% ============================================================
% causal_sessions.tex
% STATUS: STUB
% Corresponds to: src/core/CausalSession.py, src/core/PhaseRotor.py
% Experiment: exp_01_inertia.py
% ============================================================

% OPEN QUESTION (from exp_01 results):
% In the current tick() implementation, higher omega causes FASTER propagation
% given the same phase gradient, because omega amplifies the phase-alignment
% weighting. True inertial behavior requires heavier particles to resist
% changes in velocity -- omega should damp the phase-alignment response,
% not amplify it. The correct formulation needs further work.
% Candidate fix: weight emission by 1/(omega * phase_alignment) so that
% high-frequency packets are less responsive to the gradient per tick.
% This is the key open theoretical question for Section 3.


% - The CausalSession: a particle as a persistent probability flux
% - The PhaseRotor: U(1) internal clock
%   - omega as instruction frequency (mass)
%   - Phase gradient as momentum
% - The hop probability distribution
%   - At rest: 1/6 to each neighbor, probability of staying put
%   - In motion: asymmetric distribution encoding velocity
%   - The temperature analogy for the stay-put probability
% - Compact support: probability strictly zero outside causal cone
% - Inertia as phase-gradient persistence (U(1) stability)
% - Mass as instruction overhead: F = ma as state-update cost

\textbf{[STUB]}
